\documentclass[11pt]{article}
\usepackage[utf8]{inputenc}
\usepackage{amsmath,amssymb,amsthm}
\usepackage[margin=1in]{geometry}
\usepackage{enumitem}
\usepackage{hyperref}

\begin{document}
% ==================== PAGE 21 ====================
\noindent \textbf{Field:} \\
A commutative ring $(F, +, \cdot)$ with unity ($1 \ne 0$) is called a \textbf{field} if every non-zero element of $F$ has a multiplicative inverse in $F$. \\
i.e., $\forall\, a \in F \setminus \{0\}$, $\exists\, a^{-1} \in F$ such that:
\[
a \cdot a^{-1} = a^{-1} \cdot a = 1
\]

\bigskip
\noindent In other words, $(F, +, \cdot)$ is a field if:
\begin{enumerate}
    \item $(F, +)$ is an abelian group
    \item $(F \setminus \{0\}, \cdot)$ is an abelian group
    \item Distributive laws hold
\end{enumerate}

\bigskip
\noindent \textbf{Examples of Fields:}
\begin{enumerate}
    \item $(\mathbb{Q}, +, \cdot)$, $(\mathbb{R}, +, \cdot)$, $(\mathbb{C}, +, \cdot)$ are fields.
    \item $(\mathbb{Z}_p, +_p, \times_p)$ is a field if and only if $p$ is prime.
\end{enumerate}

\bigskip
\noindent \textbf{Non-examples:}
\begin{enumerate}
    \item $(\mathbb{Z}, +, \cdot)$ is not a field (only $\pm 1$ have multiplicative inverses).
    \item $(\mathbb{Z}_n, +_n, \times_n)$ where $n$ is composite is not a field.
\end{enumerate}

\bigskip
\noindent \textbf{Skew-Field / Division Ring:} \\
A ring $(R, +, \cdot)$ with unity ($1 \ne 0$) is called a \textbf{division ring} (or \textbf{skew-field}) if every non-zero element has a multiplicative inverse in $R$ (multiplication need not be commutative).

\bigskip
\noindent \textbf{Theorem:} \\
Every field is an integral domain.

\medskip
\noindent \textbf{Proof:} \\
Let $F$ be a field. Then $F$ is a commutative ring with unity. \\
To show $F$ is an integral domain, we must show it has no zero divisors. \\
Let $a, b \in F$ such that:
\[
ab = 0
\]
If $a = 0$, then there is nothing to prove. \\
If $a \ne 0$, then since $F$ is a field, $a^{-1} \in F$ exists. \\
Multiplying both sides by $a^{-1}$:
\begin{align*}
a^{-1}(ab) &= a^{-1} \cdot 0 \\
(a^{-1}a)b &= 0 \\
1 \cdot b &= 0 \\
b &= 0
\end{align*}
Thus, $ab = 0 \implies a = 0$ or $b = 0$. \\
Hence, $F$ has no zero divisors, so $F$ is an integral domain.

\newpage
% ==================== PAGE 22 ====================
\noindent \textbf{Remark:} \\
The converse of the above theorem is not true in general. \\
For example, $(\mathbb{Z}, +, \cdot)$ is an integral domain but not a field.

\bigskip
\noindent \textbf{Theorem:} \\
Every finite integral domain is a field.

\medskip
\noindent \textbf{Proof:} \\
Let $D$ be a finite integral domain. \\
Let $D = \{x_1, x_2, \dots, x_n\}$ where all elements are distinct. \\
Let $a \in D$ such that $a \ne 0$. \\
Consider the set:
\[
aD = \{a x_1, a x_2, \dots, a x_n\}
\]
Since $D$ is closed under multiplication, $aD \subseteq D$. \\
Now we show all elements of $aD$ are distinct: \\
Let $a x_i = a x_j$ for $i \ne j$. \\
Since $D$ is an integral domain and $a \ne 0$, by cancellation law:
\[
x_i = x_j
\]
which is a contradiction since elements of $D$ are distinct. \\
Thus, $aD$ contains $n$ distinct elements of $D$. \\
Since $D$ has exactly $n$ elements, we have:
\[
aD = D
\]
Since $1 \in D$ (as $D$ has unity), $\exists\, x_k \in D$ such that:
\[
a x_k = 1
\]
Since $D$ is commutative:
\[
a x_k = x_k a = 1
\]
$\Rightarrow x_k = a^{-1}$. \\
Thus, every non-zero element of $D$ has a multiplicative inverse in $D$. \\
Hence, $D$ is a field.

\bigskip
\noindent \textbf{Corollary:} \\
$(\mathbb{Z}_p, +_p, \times_p)$ is a field if and only if $p$ is a prime.

\newpage
% ==================== PAGE 23 ====================
\noindent \textbf{Characteristic of a Ring:} \\
Let $(R, +, \cdot)$ be a ring. If there exists a least positive integer $n$ such that:
\[
n \cdot a = \underbrace{a + a + \dots + a}_{n \text{ times}} = 0 \quad \forall\, a \in R
\]
then $n$ is called the \textbf{characteristic} of the ring $R$, denoted as $\text{char}(R) = n$. \\
If no such positive integer exists, then $R$ is said to be of \textbf{characteristic zero} (or infinite).

\bigskip
\noindent \textbf{Examples:}
\begin{enumerate}
    \item $\text{char}(\mathbb{Z}) = 0$, $\text{char}(\mathbb{Q}) = 0$, $\text{char}(\mathbb{R}) = 0$, $\text{char}(\mathbb{C}) = 0$.
    \item $\text{char}(\mathbb{Z}_n) = n$.
    \item $\text{char}(\mathbb{Z}_p) = p$.
\end{enumerate}

\bigskip
\noindent \textbf{Theorem:} \\
Let $R$ be a ring with unity $1$.
\begin{enumerate}
    \item If the order of $1$ under addition is finite, say $n$, then $\text{char}(R) = n$.
    \item If the order of $1$ under addition is infinite, then $\text{char}(R) = 0$.
\end{enumerate}

\medskip
\noindent \textbf{Proof:} \\
1) Let $O(1) = n$ in $(R, +)$. \\
Then $n \cdot 1 = 0$, and $n$ is the least positive integer with this property. \\
For any $a \in R$:
\[
n \cdot a = a + a + \dots + a = (1 + 1 + \dots + 1)a = (n \cdot 1)a = 0 \cdot a = 0
\]
Thus, $n \cdot a = 0$ for all $a \in R$. \\
Since $n$ is the least positive integer for which $n \cdot 1 = 0$, it is also the least positive integer for all elements. \\
$\therefore \text{char}(R) = n$.

\medskip
\noindent 2) If $O(1)$ is infinite, then there is no positive integer $n$ such that $n \cdot 1 = 0$. \\
Hence, no such $n$ exists for all elements, so $\text{char}(R) = 0$.

\bigskip
\noindent \textbf{Theorem:} \\
The characteristic of an integral domain (or a field) is either $0$ or a prime number.

\medskip
\noindent \textbf{Proof:} \\
Let $D$ be an integral domain. If $\text{char}(D) = 0$, then we are done. \\
Suppose $\text{char}(D) = n > 0$. We need to show that $n$ is a prime. \\
If possible, let $n$ be composite, so $n = n_1 n_2$, where $1 < n_1, n_2 < n$. \\
Then:
\begin{align*}
n \cdot 1 = 0 &\implies (n_1 n_2) \cdot 1 = 0 \\
&\implies (n_1 \cdot 1)(n_2 \cdot 1) = 0
\end{align*}
Since $D$ is an integral domain (no zero divisors):
\[
n_1 \cdot 1 = 0 \quad \text{or} \quad n_2 \cdot 1 = 0
\]
This contradicts the fact that $n$ is the least positive integer such that $n \cdot 1 = 0$ (since $n_1, n_2 < n$). \\
Hence, $n$ cannot be composite. \\
$\therefore n$ must be a prime number.

\end{document}
